% benötigt im Präambel: \usepackage{booktabs}
% Abb. 5.24 als gesetzte Tabelle statt Bild (Inhalt aus dem eingebetteten Objekt der 5. Auflage). Die figure-Umgebung
% hält die Nummer „Abb. 5.24“. Die beiden Teiltabellen des Originals hatten dieselben Berichtszeilen; sie stehen jetzt
% als zwei Zuordnungsspalten nebeneinander; das Pluszeichen steht vor jedem weiteren Summanden. Fett: Umsatzstufen
% wie im Original.
\begin{figure}[htbp]
  \centering
  \footnotesize
  \begin{tabular}{@{}l@{\hspace{2em}}l@{\hspace{2em}}l@{}}
    \toprule
     & \multicolumn{2}{l@{}}{Kennzahlenzuordnung} \\
    \cmidrule(l{0pt}){2-3}
    Berichtszeile & eine Kennzahl je Zeile & Summe mehrerer Kennzahlen \\
    \midrule
    \textbf{Bruttoumsatz}      & Bruttoumsatz      & Bruttoumsatz \\
    Rabatte                    & Rabatte           & Kundenrabatte \\
                               &                   & \llap{$+$\,}Aktionsrabatte \\
                               &                   & \llap{$+$\,}Sonderrabatte \\
    \textbf{Nettoumsatz}       & Nettoumsatz       & Nettoumsatz \\
    Skonti                     & Skonti            & Skonti \\
    \textbf{NettoNetto-Umsatz} & NettoNetto-Umsatz & NettoNetto-Umsatz \\
    \bottomrule
  \end{tabular}
  \caption{Kennzahlenzuordnung zu Berichtszeilen}
  \label{fig:Kennzahlenzuordnung_zu_Berichtszeilen}
\end{figure}
